% GENERATED FILE. Edit canonical lessons, then run npm run generate:books.
% canonical-source-hash: fnv1a64:61b501b9ba275060
% canonical-lessons: IT-C515-pattinare, IT-C515-scarpone, IT-C515-pista, IT-C515-sciatore, IT-C515-seggiovia

\chapter{Ghiaccio e neve}
\label{ch:it-ghiaccio-e-neve}
\hlchaptermodality{515}

\begin{chapteropening}
\textbf{By the end of this chapter:} \emph{I can talk about winter sports.}

Complete the last lesson of chapter 515: I can talk about winter sports.
\end{chapteropening}

\section[pattinare]{pattinare — to skate}

\label{lesson:IT-C515-pattinare}



\begin{quote}
Before the new one: say the Italian for to dive under, then the Italian for diving board.
\end{quote}

\subsection*{You'll want to know: pattinare}
\textbf{pattinare} — \textquotedblleft{}to skate\textquotedblright{}.

\begin{quote}
D'inverno pattiniamo sul ghiaccio in piazza.
\end{quote}

\begin{quote}
\emph{In winter we skate on the ice in the square.}
\end{quote}

Skates are \emph{i pattini}.

\begin{cousinweb}
\emph{pattinare} is built on \emph{pattino} \textquotedblleft{}skate\textquotedblright{}; the origin of \emph{pattino} is not settled, and it may come from an old word \emph{patta} \textquotedblleft{}paw\textquotedblright{}.
\end{cousinweb}

\begin{grammarlens}[title={What you've built}]
One more word for telling, explaining and describing at B1.
\end{grammarlens}

\subsection*{Your turn}
\begin{itemize}
\raggedright
  \item \emph{Say it:} \emph{pattinare}
  \item \emph{Say it:} \emph{pattinare}, once more
  \item \emph{Say it:} a sentence of your own with \emph{pattinare}
  \item \emph{Recall:} say the Italian for to dive under, then the Italian for diving board, then say \emph{pattinare} again
\end{itemize}

\begin{tcolorbox}[breakable,colback=teal!4,colframe=teal!35!black,title={Before you move on}]
What does pattinare mean? (\textquotedblleft{}To skate\textquotedblright{}.) Say it once more.
\end{tcolorbox}
\section[lo scarpone]{lo scarpone — boot (for skiing or hiking)}

\label{lesson:IT-C515-scarpone}



\begin{quote}
Before the new one: say the Italian for diving board, then the Italian for to skate.
\end{quote}

\subsection*{You'll want to know: lo scarpone}
\textbf{lo scarpone} — \textquotedblleft{}boot (for skiing or hiking)\textquotedblright{}.

\begin{quote}
Gli scarponi da sci sono pesanti.
\end{quote}

\begin{quote}
\emph{Ski boots are heavy.}
\end{quote}

A big \emph{scarpa}: \emph{-one} makes things bigger.

\begin{cousinweb}
\emph{scarpone} is built on \emph{scarpa} \textquotedblleft{}shoe\textquotedblright{} with \emph{-one} \textquotedblleft{}big\textquotedblright{}; \emph{scarpa} is usually traced to a Germanic word, and its origin is not settled.
\end{cousinweb}

\begin{grammarlens}[title={What you've built}]
One more word for telling, explaining and describing at B1.
\end{grammarlens}

\subsection*{Your turn}
\begin{itemize}
\raggedright
  \item \emph{Say it:} \emph{lo scarpone}
  \item \emph{Say it:} \emph{lo scarpone}, once more
  \item \emph{Say it:} a sentence of your own with \emph{lo scarpone}
  \item \emph{Recall:} say the Italian for diving board, then the Italian for to skate, then say \emph{lo scarpone} again
\end{itemize}

\begin{tcolorbox}[breakable,colback=teal!4,colframe=teal!35!black,title={Before you move on}]
What does lo scarpone mean? (\textquotedblleft{}Boot (for skiing or hiking)\textquotedblright{}.) Say it once more.
\end{tcolorbox}
\section[la pista]{la pista — slope, track, runway}

\label{lesson:IT-C515-pista}



\begin{quote}
Before the new one: say the Italian for to skate, then the Italian for boot.
\end{quote}

\subsection*{You'll want to know: la pista}
\textbf{la pista} — \textquotedblleft{}slope, track, runway\textquotedblright{}.

\begin{quote}
Oggi le piste sono piene di neve fresca.
\end{quote}

\begin{quote}
\emph{Today the slopes are full of fresh snow.}
\end{quote}

\emph{La pista da sci}, \emph{la pista ciclabile}, \emph{la pista da ballo} — the dance floor.

\begin{cousinweb}
\emph{pista} is usually traced to \emph{pestare} \textquotedblleft{}to tread\textquotedblright{}, from Late Latin \emph{pistāre} \textquotedblleft{}to pound\textquotedblright{}, built on Latin \emph{pinsere} \textquotedblleft{}to crush\textquotedblright{} (Lewis \& Short): a track is ground trodden flat.
\end{cousinweb}

\begin{grammarlens}[title={What you've built}]
One more word for telling, explaining and describing at B1.
\end{grammarlens}

\subsection*{Your turn}
\begin{itemize}
\raggedright
  \item \emph{Say it:} \emph{la pista}
  \item \emph{Say it:} \emph{la pista}, once more
  \item \emph{Say it:} a sentence of your own with \emph{la pista}
  \item \emph{Recall:} say the Italian for to skate, then the Italian for boot, then say \emph{la pista} again
\end{itemize}

\begin{tcolorbox}[breakable,colback=teal!4,colframe=teal!35!black,title={Before you move on}]
What does la pista mean? (\textquotedblleft{}Slope, track, runway\textquotedblright{}.) Say it once more.
\end{tcolorbox}
\section[lo sciatore]{lo sciatore — skier}

\label{lesson:IT-C515-sciatore}



\begin{quote}
Before the new one: say the Italian for boot, then the Italian for slope.
\end{quote}

\subsection*{You'll want to know: lo sciatore}
\textbf{lo sciatore} — \textquotedblleft{}skier\textquotedblright{}.

\begin{quote}
Ci sono molti sciatori sulla pista.
\end{quote}

\begin{quote}
\emph{There are lots of skiers on the slope.}
\end{quote}

A woman is \emph{la sciatrice}.

\begin{cousinweb}
\emph{sciatore} is built on \emph{sci} \textquotedblleft{}ski\textquotedblright{}, which was borrowed from Norwegian \emph{ski}, from Old Norse \emph{skíð} \textquotedblleft{}a stick of wood, a snowshoe\textquotedblright{} (EtymDB, under French \emph{ski}).
\end{cousinweb}

\begin{grammarlens}[title={What you've built}]
One more word for telling, explaining and describing at B1.
\end{grammarlens}

\subsection*{Your turn}
\begin{itemize}
\raggedright
  \item \emph{Say it:} \emph{lo sciatore}
  \item \emph{Say it:} \emph{lo sciatore}, once more
  \item \emph{Say it:} a sentence of your own with \emph{lo sciatore}
  \item \emph{Recall:} say the Italian for boot, then the Italian for slope, then say \emph{lo sciatore} again
\end{itemize}

\begin{tcolorbox}[breakable,colback=teal!4,colframe=teal!35!black,title={Before you move on}]
What does lo sciatore mean? (\textquotedblleft{}Skier\textquotedblright{}.) Say it once more.
\end{tcolorbox}
\section[la seggiovia]{la seggiovia — chairlift}

\label{lesson:IT-C515-seggiovia}



\begin{quote}
Before the new one: say the Italian for slope, then the Italian for skier.
\end{quote}

\subsection*{You'll want to know: la seggiovia}
\textbf{la seggiovia} — \textquotedblleft{}chairlift\textquotedblright{}.

\begin{quote}
Saliamo in seggiovia fino in cima.
\end{quote}

\begin{quote}
\emph{We go up to the top by chairlift.}
\end{quote}

Built from \emph{seggio} \textquotedblleft{}seat\textquotedblright{} and \emph{via} \textquotedblleft{}way\textquotedblright{}.

\begin{cousinweb}
\emph{seggiovia} is a modern word built from \emph{seggio} \textquotedblleft{}seat\textquotedblright{}, from Latin \emph{sedēre} \textquotedblleft{}to sit\textquotedblright{} (Lewis \& Short), and \emph{via} \textquotedblleft{}way\textquotedblright{}.
\end{cousinweb}

\begin{grammarlens}[title={What you've built}]
One more word for telling, explaining and describing at B1.
\end{grammarlens}

\subsection*{Your turn}
\begin{itemize}
\raggedright
  \item \emph{Say it:} \emph{la seggiovia}
  \item \emph{Say it:} \emph{la seggiovia}, once more
  \item \emph{Say it:} a sentence of your own with \emph{la seggiovia}
  \item \emph{Recall:} say the Italian for slope, then the Italian for skier, then say \emph{la seggiovia} again
\end{itemize}

\begin{tcolorbox}[breakable,colback=teal!4,colframe=teal!35!black,title={Before you move on}]
What does la seggiovia mean? (\textquotedblleft{}Chairlift\textquotedblright{}.) Say it once more.
\end{tcolorbox}
